\chapter{Introduction to Environment and Sustainable Development}
\label{chap:intro-sustainability}

\section{Overview and Course Learning Outcomes}
Environmental Studies is an inherently multidisciplinary academic endeavor that investigates the intricate physical, chemical, biological, social, and economic processes that sustain life on Earth. As anthropogenic activities exert unprecedented stresses on planetary boundaries, understanding ecological systems and transitioning toward sustainable development paradigms has become an urgent mandate for contemporary engineers, scientists, and decision-makers.

\begin{important}[Course Learning Outcomes for Unit 1]
Upon completing this chapter, the student will be capable of:
\begin{enumerate}[leftmargin=1.5em]
    \item Defining the constituent spheres of the Earth system (Atmosphere, Hydrosphere, Lithosphere, Biosphere) and analyzing their continuous biogeochemical mass-energy exchanges.
    \item Explaining the multidisciplinary nature, professional scope, and societal necessity of environmental studies.
    \item Formulating and quantitatively evaluating the Ehrlich-Holdren IPAT model ($I = P \times A \times T$) and carrying capacity constraints.
    \item Differentiating between weak and strong sustainability paradigms under neoclassical versus ecological economics.
    \item Outlining the 17 United Nations Sustainable Development Goals (SDGs), their key indicators, interlinkages, and monitoring frameworks.
\end{enumerate}
\end{important}

\section{The Environment: Constituents, Structure and Earth Spheres}
\begin{definition}[Environment]
The environment is defined as the totality of surrounding external physical, chemical, biological, and socio-economic factors that act upon an organism or an ecological community during its lifespan, ultimately dictating its morphology, physiological function, distribution, and survival.
\end{definition}

The planetary environment functions as an open thermodynamic system receiving low-entropy solar shortwave radiation ($\sim 1361\text{ W/m}^2$ solar constant) and radiating degraded high-entropy longwave infrared radiation back into deep space. Structurally, the Earth system is partitioned into four major interacting spheres:

\subsection{The Atmosphere}
The atmosphere is a dynamic gaseous envelope of approximately $5.15 \times 10^{18}\text{ kg}$ held gravitationally around the globe. Chemically, clean, dry air at sea level comprises $78.08\%$ Nitrogen ($N_2$), $20.95\%$ Oxygen ($O_2$), $0.93\%$ Argon ($Ar$), $0.042\%$ ($420\text{ ppm}$) Carbon Dioxide ($CO_2$), along with trace concentrations of neon, helium, methane, krypton, and nitrous oxide. Water vapor concentration varies spatio-temporally from $<0.1\%$ in polar regions to $4.0\%$ in humid tropical marine boundary layers.

Thermally, the atmosphere exhibits a vertical stratification characterized by alternating lapse rate gradients:
\begin{enumerate}[leftmargin=1.5em]
    \item \textbf{Troposphere ($0\text{--}12\text{ km}$):} Houses more than $80\%$ of atmospheric mass and virtually all atmospheric moisture and weather phenomena. Temperature decreases with altitude at the Environmental Lapse Rate (ELR $\approx 6.5^\circ\text{C/km}$) until reaching the tropopause ($-56^\circ\text{C}$).
    \item \textbf{Stratosphere ($12\text{--}50\text{ km}$):} Houses the protective ozone layer ($O_3$), primarily concentrated between $20\text{--}30\text{ km}$. Absorption of solar ultraviolet radiation ($\text{UV-B}$ and $\text{UV-C}$) produces a thermal inversion where temperature increases with altitude, reaching $\approx 0^\circ\text{C}$ at the stratopause.
    \item \textbf{Mesosphere ($50\text{--}85\text{ km}$):} Characterized by falling temperatures down to $-90^\circ\text{C}$ (the coldest region of the Earth system); site of meteoric particle ablation.
    \item \textbf{Thermosphere ($85\text{--}600\text{ km}$):} Extreme solar UV and X-ray absorption induces intense ionization (forming the ionosphere) and elevates kinetic temperatures above $1200^\circ\text{C}$, though molecular density is negligible.
\end{enumerate}

\subsection{The Hydrosphere, Lithosphere, and Biosphere}
\begin{itemize}[leftmargin=1.5em]
    \item \textbf{Hydrosphere:} Represents the total aqueous reservoir of the planet, estimated at $1.386 \times 10^9\text{ km}^3$. Of this volume, $97.5\%$ consists of saline oceanic waters. The remaining $2.5\%$ comprises global freshwater, whereof $68.7\%$ is sequestered in polar ice sheets and glaciers, $30.1\%$ resides as groundwater, and only $1.2\%$ exists as accessible surface freshwater (rivers, lakes, and freshwater wetlands).
    \item \textbf{Lithosphere:} The rigid outermost silicate shell encompassing the Earth's crust and uppermost mantle. Continental crust consists predominantly of felsic granitic rocks ($2.7\text{ g/cm}^3$, thickness $30\text{--}50\text{ km}$), whereas oceanic crust is composed of denser mafic basalts ($3.0\text{ g/cm}^3$, thickness $5\text{--}10\text{ km}$). Weathering of lithospheric rocks provides essential mineral nutrients (phosphorus, potassium, calcium) driving biological productivity.
    \item \textbf{Biosphere:} The thin, life-supporting veneer encompassing the lower atmosphere, the entire hydrosphere, and the uppermost weathered layer of the lithosphere (pedosphere).
\end{itemize}

\section{Environmental Studies as a Multidisciplinary Field}
Environmental problems cannot be adequately conceptualized or solved through the lens of a single traditional discipline. An effective solution to industrial air pollution, for instance, requires:
\begin{itemize}[leftmargin=1.5em]
    \item \textbf{Atmospheric Physics and Chemistry:} To model turbulent dispersion, plume rise, and photochemical smog reactions.
    \item \textbf{Ecology and Toxicology:} To quantify physiological damage thresholds in flora, fauna, and human respiratory systems.
    \item \textbf{Chemical and Mechanical Engineering:} To design scrubbing equipment, electrostatic precipitators, and catalytic converters.
    \item \textbf{Economics and Jurisprudence:} To formulate Pigouvian carbon taxes, cap-and-trade markets, and statutory emission standards.
\end{itemize}

\section{Quantitative Sustainability and Carrying Capacity}
\begin{theorem}[Brundtland Principle of Intergenerational Equity]
Sustainable development is development that meets the needs of the present without compromising the ability of future generations to meet their own needs. It establishes two fundamental concepts: (1) the concept of `needs', particularly the essential needs of the world's poor, and (2) the limitations imposed by the state of technology and social organization on the environment's ability to meet present and future needs.
\end{theorem}

\begin{keyequation}[The Ehrlich-Holdren IPAT Model]
Anthropogenic environmental impact ($I$) is driven by three multiplicative macro-drivers:
\begin{equation}
I = P \times A \times T
\end{equation}
where:
\begin{itemize}[leftmargin=1.5em]
    \item $P = \text{Population size (number of human individuals)}$.
    \item $A = \text{Affluence (per-capita consumption of goods, measured as GDP/person)}$.
    \item $T = \text{Technology intensity (environmental degradation or greenhouse gas emissions per unit consumption)}$.
\end{itemize}
\end{keyequation}

\begin{example}[IPAT Decarbonization Requirement]
A country has an initial population $P_0 = 100\text{ million}$, per-capita GDP $A_0 = \$10{,}000$, and technology intensity $T_0 = 0.5\text{ kg } CO_2/\$$. Over a 20-year horizon, $P$ increases by $20\%$ ($P_1 = 1.2 P_0$) and $A$ triples ($A_1 = 3.0 A_0$). If national policy mandates that total emissions $I$ must decrease by $40\%$ ($I_1 = 0.60 I_0$), determine the required target technology intensity $T_1$ and its percentage reduction.
\end{example}

\begin{solutionbox}
Initial impact: $I_0 = P_0 \times A_0 \times T_0$. \\
Future impact constraint: $I_1 = 0.60 I_0 = 0.60 (P_0 A_0 T_0)$. \\
Expressing future impact via IPAT:
\begin{equation*}
I_1 = P_1 \times A_1 \times T_1 = (1.2 P_0) \times (3.0 A_0) \times T_1 = 3.6 (P_0 A_0) T_1.
\end{equation*}
Equating the two expressions for $I_1$:
\begin{equation*}
3.6 (P_0 A_0) T_1 = 0.60 (P_0 A_0 T_0) \implies T_1 = \frac{0.60}{3.6} T_0 = \frac{1}{6} T_0 \approx 0.1667 T_0.
\end{equation*}
Substituting $T_0 = 0.5\text{ kg } CO_2/\$$:
\begin{equation*}
T_1 = \frac{0.5}{6} = 0.0833\text{ kg } CO_2/\$.
\end{equation*}
The required technological reduction is:
\begin{equation*}
\Delta T\% = \left(1 - \frac{T_1}{T_0}\right) \times 100\% = \left(1 - \frac{1}{6}\right) \times 100\% = 83.33\%.
\end{equation*}
The carbon intensity per dollar of GDP must decline by $83.33\%$ to offset population expansion and economic affluence growth.
\end{solutionbox}

\section{The United Nations Sustainable Development Goals (SDGs)}
Adopted by all 193 UN Member States in 2015 as part of the 2030 Agenda for Sustainable Development, the 17 SDGs encompass 169 actionable targets and 232 unique monitoring indicators:
\begin{itemize}[leftmargin=1.5em]
    \item \textbf{Basic Human Needs (SDGs 1--5):} Eradicating extreme poverty, ending hunger, ensuring universal healthcare, quality education, and gender equality.
    \item \textbf{Planetary Biophysical Boundaries (SDGs 6, 13, 14, 15):} SDG 6 (Clean Water and Sanitation), SDG 13 (Climate Action), SDG 14 (Life Below Water), and SDG 15 (Life on Land).
    \item \textbf{Economic Transformation (SDGs 7, 8, 9, 11, 12):} SDG 7 (Affordable and Clean Energy), SDG 11 (Sustainable Cities and Communities), and SDG 12 (Responsible Consumption and Production).
    \item \textbf{Governance and Partnerships (SDGs 16--17):} Peace, justice, strong institutions, and global multi-stakeholder partnerships.
\end{itemize}

\begin{quickrevision}[Chapter 1 Summary]
\begin{itemize}[leftmargin=1.2em]
    \item The environment comprises four interacting spheres: Atmosphere, Hydrosphere, Lithosphere, and Biosphere.
    \item The troposphere contains $80\%$ of air mass; the stratosphere contains the ozone layer ($20\text{--}30\text{ km}$).
    \item Freshwater comprises only $2.5\%$ of global water, with $68.7\%$ trapped in ice caps and $30.1\%$ in groundwater.
    \item The IPAT identity ($I = P \times A \times T$) links environmental impact to population, affluence, and technology.
    \item The 17 SDGs balance economic growth, social inclusion, and biophysical environmental limits by 2030.
\end{itemize}
\end{quickrevision}
